\documentclass[12pt, english]{article}
\usepackage[margin = 1.5in]{geometry}
%%%%%%%%%%%%%%%%%%%%%%%%%%%%%%%%%%%%%%%%%%%%%%%%%%%%%%%%%%%%%%%%%%%%%%%%%%%%%%%%%% PACKAGES
\usepackage[T1]{fontenc}
\usepackage[english]{babel}
\usepackage[shortlabels]{enumitem}
\usepackage{amsmath, amssymb,url,xspace, amsmath, floatrow, hyperref, graphicx}
\usepackage{subfiles}
\usepackage[ruled,vlined]{algorithm2e}
\usepackage{diagbox}
\usepackage{color, colortbl}
\definecolor{LightCyan}{rgb}{0.88,1,1}

\usepackage{pdfpages}
\usepackage{hyperref}
\hypersetup{
	colorlinks=true,
	linkcolor=blue,
	filecolor=magenta,      
	urlcolor=cyan,
	citecolor=red
}


\usepackage{subcaption}

%%%%%%%%%%%%%%%%%%%%%%%%%%%%%%%%%%%%%%%%%%%%%%%%%%%%%%%%%%%%%%%%%%%%%%%%%%%%%%%%%% CUSTOM COMMANDS
\newcommand{\argmin}{\text{arg} \min}
\newcommand{\sgn}[1]{\text{sgn} \left ( #1 \right )}
\newcommand{\p}{\mathbf{p}}
\newcommand{\W}{\mathbf{W}}
\newcommand{\U}{\mathbf{U}}
\newcommand{\D}{\mathbf{D}}
\newcommand{\V}{\mathbf{V}}
\newcommand{\X}{\mathbf{X}}
\newcommand{\y}{\mathbf{y}}
\newcommand{\z}{\mathbf{z}}
\newcommand{\bbeta}{\boldsymbol \beta}
\newcommand{\bgamma}{\boldsymbol \gamma}
\newcommand{\old}{\text{old}}
\newcommand{\new}{\text{new}}
\newcommand{\ridge}{\text{ridge}}
\newcommand{\lasso}{\text{lasso}}
\newcommand{\RgPCR}{\text{RgPCR}}
\newcommand{\RgPCC}{\text{RgPCC}}
\newcommand{\ols}{\text{ols}}
\renewcommand{\exp}[1]{\text{exp} \left ( #1 \right )}

\newcommand{\gammaone}{\bgamma_1 = (25, 0, \ldots , 0)}
\newcommand{\gammatwo}{\bgamma_2 = (20, 10, 10, 0, \ldots , 0)}
\newcommand{\gammathree}{\bgamma_3 = (15, 10, 5, 5, 3, 0, \ldots , 0)}

\renewcommand{\tilde}[1]{\widetilde{#1}}
\renewcommand{\hat}[1]{\widehat{#1}}

\def\meta#1{$\langle${\it #1}$\rangle$}
\newcommand{\myast}{($\star$)\ }
\makeatletter
       \def\ps@copyright{\ps@empty
    \def\@oddfoot{\hfil\small RTG Mini-Conference}}
\makeatother

\newcommand{\BibTeX}{{\scshape Bib}\kern-.08em\TeX}
\newcommand{\T}{\S\kern .15em\relax }
\newcommand{\AMS}{$\mathcal{A}$\kern-.1667em\lower.5ex\hbox
        {$\mathcal{M}$}\kern-.125em$\mathcal{S}$}
\newcommand{\resp}{\emph{resp}.\xspace}

%%%%%%%%%%%%%%%%%%%%%%%%%%%%%%%%%%%%%%%%%%%%%%%%%%%%%%%%%%%%%%%%%%%%%%%%%%%%%%%%%% SETTINGS
\tolerance 400
\pretolerance 200

%%%%%%%%%%%%%%%%%%%%%%%%%%%%%%%%%%%%%%%%%%%%%%%%%%%%%%%%%%%%%%%%%%%%%%%%%%%%%%%%%% NOTATION AND LIST OF SYMBOLS STUFF


\usepackage[sort = def, nonumberlist]{glossaries}
\renewcommand*{\glstextformat}[1]{\textcolor{black}{#1}}
\newglossary{symbols}{sym}{sbl}{List of Abbreviations and Symbols}
\makeglossary



% data stuff
\newglossaryentry{X}{type=symbols,name={\ensuremath{\X}},
	description={a design matrix, usually with dimensions $N \times p$}}

\newglossaryentry{y}{type=symbols,name={\ensuremath{\y}},
	description={a response vector of the binary classes}}

\newglossaryentry{U}{type=symbols,name={\ensuremath{\U}},
	description={of the svd of $\X = \U \D \V^T$}}

\newglossaryentry{D}{type=symbols,name={\ensuremath{\D}},
	description={of the svd of $\X = \U \D \V^T$}}

\newglossaryentry{V}{type=symbols,name={\ensuremath{\V}},
	description={of the svd of $\X = \U \D \V^T$}}

\newglossaryentry{W}{type=symbols,name={\ensuremath{\W}},
	description={a weight matrix}}

\newglossaryentry{z}{type=symbols,name={\ensuremath{\z}},
	description={a psuedo response based on the weight $\W$}}

% coef stuff


%%%%%%%%%%%%%%%%%%%%%%%%%%%%%%%%%%%%%%%%%%%%%%%%%%%%%%%%%%%%%%%%%%%%%%%%%%%%%%%%%% TITLE
%\title{Response-guided Principal Component Classification}
%\date {Nov 2020}
%\author{Duncan Bennett \\ Adviser : Hao Helen Zhang}


%%%%%%%%%%%%%%%%%%%%%%%%%%%%%%%%%%%%%%%%%%%%%%%%%%%%%%%%%%%%%%%%%%%%%%%%%%%%%%%%%% BEGIN DOCUMENT
\begin{document}
\def\smfbyname{}

\includepdf[page={1}]{titlepage.pdf} 

place holder for signature page
\newpage

% the acknowledgements page

\hspace{0pt}
\vfill
\begin{center}
	Thank you to my adviser Helen Zhang, to my wife Samantha Kao and to my family. This would not have been possible without you.
\end{center}
\vfill
\hspace{0pt}

\newpage

\tableofcontents
\newpage

\printglossaries
\newpage

\listoffigures
\newpage

%\listoftables
%\newpage

\begin{abstract}
\addcontentsline{toc}{section}{Abstract}
	Response-guided principal component regression (RgPCR) is a generalization of ridge regression. This method improves upon principal component regression (PCR) by taking the response values into account during variable selection. In this paper, we will modify RgPCR for binary classification problems using ideas from logistic regression. This technique is called \textit{response-guided principal component classification (RgPCC)}.
\end{abstract}
\newpage


\subfile{intro} 


\subfile{literature-review} 


\subfile{RgPCC}


\subfile{numerical-studies}


\subfile{further-research}

%\newgeometry{top=0.75in,bottom=0.75in,right=0.75in,left=0.75in}

\section{Figures}
\subfile{figures}
%%%%%%%%%%%%%%%%%%%%%%%%%%%%%%%%%%%%%%%%%%%%%%%%%
%%%%%%%%%%%%%%%%%% BIBLIOGRAPHY %%%%%%%%%%%%%%%%%
%%%%%%%%%%%%%%%%%%%%%%%%%%%%%%%%%%%%%%%%%%%%%%%%%

\addcontentsline{toc}{section}{References}
%\def\refname{Bibliography}
\begin{thebibliography}{99}

\bibitem {elemstatlearn}
 {\sc T. \ Hastie, \ R. \ Tibshirani, \ and \ J. \ Friedman (2009)} ---
 {\it The Elements of Statistical Learning : Data Mining, Inference, and Prediction
 (second edition)}, New York, NY: Springer.

 \bibitem {langzou}
 {\sc W. \ Lang, \ and \ H. \ Zou (2020)} ---
 {\it A simple method to improve principal components regression},
 WILEY, \href{https://onlinelibrary.wiley.com/doi/abs/10.1002/sta4.288}{https://onlinelibrary.wiley.com/doi/abs/10.1002/sta4.288}

 \bibitem {wangleng}
 {\sc H. \ Wang \ and \ C. \ Leng (2007)} ---
 {\it Unified LASSO Estimation by Least Squares Approximation},
 Journal of the American Statistical Association, 102:479, 1039-1048, DOI; 10.1198/016214507000000509, \href{https://doi.org/10.1198/016214507000000509}{https://doi.org/10.1198/016214507000000509}

\end{thebibliography}


\end{document}
















% unused figures

%=========================================================== old results tables

%\begin{figure}[H]
%	\begin{tabular}{c|c|c|c|c}
%		$N = $ sample size & $\bgamma_i$ & $\lambda$ & $MSE(\hat{\p}^{\log})$ & $MSE(\hat{\p}^{\RgPCC})$ \\ \hline
%		100 & $\bgamma_1$ & 43 & 0.04490677 & 0.00129034\\
%		200 & $\bgamma_1$ & 114 & 0.02896408 & 0.0006097103\\
%		100 & $\bgamma_2$ & 19 & 0.04305285 & 0.006591401 \\
%		200 & $\bgamma_2$ & 38 & 0.0281505 & 0.005462017 \\
%		100 & $\bgamma_3$ & 26 & 0.04519566  & 0.008713318  \\
%		200 & $\bgamma_3$ & 54 & 0.02590808 & 0.005361949
%	\end{tabular}
%\caption{Testing error for dimension $p = 12$}
%\end{figure}


% \begin{figure}[H]
% 	\begin{tabular}{c|c|c|c|c}
% 		sample size & $\bgamma_i$ & $\lambda$ & $MSE(\hat{\p}^{\log})$ & $MSE(\hat{\p}^{\RgPCC})$ \\ \hline
% 		200 & $\bgamma_1$ & 43 & 0.04490677 & 0.00129034\\
% 		500 & $\bgamma_1$ & 114 & 0.02896408 & 0.0006097103\\
% 		200 & $\bgamma_2$ & 19 & 0.04305285 & 0.006591401 \\
% 		500 & $\bgamma_2$ & 38 & 0.0281505 & 0.005462017 \\
% 		200 & $\bgamma_3$ & 26 & 0.04519566  & 0.008713318  \\
% 		500 & $\bgamma_3$ & 54 & 0.02590808 & 0.005361949
% 	\end{tabular}
% \caption{Testing error for dimension $p = 100$ (to be updated)}
% \end{figure}


